\Needspace{16\baselineskip}
\SourceLeaf{104}%

\begin{center}
REMARKS ON THE INVERSION OF THE ELLIPTIC INTEGRAL OF THE FIRST KIND.
\end{center}

\noindent\hspace*{\parindent}This note, probably datable from the diary entry of 6\ May 1800, »Theoriam quantitatum transcendentium $\int\frac{dx}{\BareRoot{(1-a xx)(1-\beta xx)}}$ ad summam universalitatem perduximus«, is among the most interesting derivations that \textsc{Gauss} carried out in the field of elliptic functions. In its formulation of the problem and its notation, it follows the investigations »Elegantiores integralis $\int\frac{dx}{\BareRoot{(1-x^4)}}$ proprietates« (ges. Werke III p.\,404 ff.), which belong to the earliest period of his work on lemniscatic functions. On the other hand, in 1799 \textsc{Gauss} studied the inversion of the general integral of the first kind $\int\frac{du}{\BareRoot{(1+\mu\mu\sin u^2)}}$.

The letter from \textsc{Gauss} to \textsc{Bessel} printed above sheds particularly interesting light on the mathematics developed here. In it \textsc{Gauss} repeatedly emphasizes the importance of entire transcendental functions. In an elegant derivation, the function $f\varphi$, which \textsc{Jacobi} later denoted by $\sin\mathrm{am}\,\varphi$, is here represented as the quotient of two functions $P\varphi$ and $Q\varphi$, which \textsc{Gauss} unquestionably knew to be \textls{entire} transcendental functions.

It may also be noted that the functions $P, Q$ that \textsc{Gauss} introduces here are precisely those to which \textsc{Weierstrass}, building on the work of \textsc{Abel}, later gave the designations $\mathrm{Al}(\varphi)_1$ and $\mathrm{Al}(\varphi)_0$.

\hfill \textsc{Fricke}.

\bigskip

